\section{\texttt{TROWARGMAX}}
\noindent\textbf{Descriptor profile.} Vec entry 42, tile catalog opcode \texttt{0xF000}. Descriptor bundle: \texttt{B.OP + B.DIM + B.IOT}; WO/WM sites: 13; VEC beat-level sites: 12; ROM bytes: 104.\par\smallskip
\TileInstructionEncoding{figs/encoding/tile_instruction_entries/entry_42_pto_trowargmax_th0.pdf}{figs/encoding/tile_instruction_entries/entry_42_pto_trowargmax_th2.pdf}{0xF000}{42}{B.OP + B.DIM + B.IOT}
\paragraph{Bundled descriptor (PTO-ISA header).}
\begin{description}[leftmargin=0.18\linewidth,style=nextline]
\item[Bundle] \texttt{B.OP + B.DIM + B.IOT}. \texttt{B.OP.desc\_count} = 2; \texttt{B.OP.rom\_entry\_id} = 42.
\item[Assembly] 0: \texttt{B.OP} selects \texttt{TROWARGMAX}, carries element format/variant, declares \texttt{desc\_count}, and names the ROM entry. 1: \texttt{B.DIM} supplies row-axis reduction or row-broadcast bounds. 2: \texttt{B.IOT} binds architectural tile operands. Across all \texttt{B.IOT} descriptors: tile inputs \texttt{src}, mask inputs none, and outputs \texttt{dst}.
\item[Tile operands] 1 tile input(s): \texttt{src}; 1 tile output(s): \texttt{dst}. Bound by 1 architectural \texttt{B.IOT} descriptor; each \texttt{B.IOT} supplies at most two source selectors plus one destination selector. ROM scratch tiles are not counted in \texttt{B.OP.desc\_count}.
\item[Scalar GPR operands] No scalar GPR import/export descriptor is required.
\item[Metadata operands] \texttt{B.DIM} supplies row-axis reduction or row-broadcast bounds.
\item[ROM scratch operands] \texttt{tmp} are allocated by lowering or the implementation profile for the microcode body. They are not encoded in the CPU-visible PTO-ISA header.
\end{description}
\par\smallskip
{\small
\paragraph{PTO ISA description.}
Get the column index of the maximum element for each row.
\paragraph{PTO ISA operation.}
Let \texttt{R = src.GetValidRow()} and \texttt{C = src.GetValidCol()}. For \texttt{0 \textless{}= i \textless{} R}:

\[
\mathrm{dst}_{i,0} = \max_{0 \le j < C} j_{i}
\]
}\par\smallskip
\paragraph{Microcode site table.}
{\scriptsize
\begin{longtable}{@{}R{0.055\linewidth}L{0.23\linewidth}L{0.31\linewidth}L{0.22\linewidth}@{}}
\toprule
Seq & Micro instruction & WO[63:32] fields & WM[31:0] fields \\
\midrule
0 & \texttt{u\_vec.vbr} & \texttt{opc=0x0B, x=0, fl=alu, seq=0, kind=c} & \texttt{entry=42, cnt=0, res=vec, var=0} \\
1 & \texttt{u\_vec.vbr} & \texttt{opc=0x0B, x=0, fl=alu, seq=1, kind=c} & \texttt{entry=42, cnt=0, res=vec, var=0} \\
2 & \texttt{u\_vec.vlds} & \texttt{opc=0x17, x=0, fl=mem, seq=2, kind=b} & \texttt{entry=42, cnt=2, res=vec, var=0} \\
3 & \texttt{u\_vec.vcmax} & \texttt{opc=0x0D, x=0, fl=alu, seq=3, kind=b} & \texttt{entry=42, cnt=2, res=vec, var=0} \\
4 & \texttt{u\_vec.vdintlv} & \texttt{opc=0x11, x=0, fl=alu, seq=4, kind=u} & \texttt{entry=42, cnt=1, res=vec, var=0} \\
5 & \texttt{u\_vec.vbr} & \texttt{opc=0x0B, x=0, fl=alu, seq=5, kind=c} & \texttt{entry=42, cnt=0, res=vec, var=0} \\
6 & \texttt{u\_vec.vbitcast} & \texttt{opc=0x0A, x=0, fl=alu, seq=6, kind=u} & \texttt{entry=42, cnt=1, res=vec, var=0} \\
7 & \texttt{u\_vec.vadds} & \texttt{opc=0x08, x=0, fl=alu, seq=7, kind=b} & \texttt{entry=42, cnt=2, res=vec, var=0} \\
8 & \texttt{u\_vec.vcmp} & \texttt{opc=0x0F, x=0, fl=alu, seq=8, kind=b} & \texttt{entry=42, cnt=2, res=vec, var=0} \\
9 & \texttt{u\_vec.vsel} & \texttt{opc=0x24, x=0, fl=alu, seq=9, kind=t} & \texttt{entry=42, cnt=3, res=vec, var=0} \\
10 & \texttt{u\_vec.pbitcast} & \texttt{opc=0x01, x=0, fl=setup, seq=10, kind=u} & \texttt{entry=42, cnt=1, res=vec, var=0} \\
11 & \texttt{u\_vec.vsel} & \texttt{opc=0x24, x=0, fl=alu, seq=11, kind=t} & \texttt{entry=42, cnt=3, res=vec, var=0} \\
12 & \texttt{u\_vec.vsts} & \texttt{opc=0x2A, x=0, fl=mem, seq=12, kind=b} & \texttt{entry=42, cnt=2, res=vec, var=0} \\
\bottomrule
\end{longtable}
}
\paragraph{ROM body DSL.}
\begin{lstlisting}[style=ptocode]
def rom_trowargmax(src, tmp, dst):
    src_dtype = src.element_type
    idx_dtype = dst.element_type
    lanes = pto.get_lanes(src_dtype)
    valid_rows, valid_cols = src.valid_shape
    elem_bytes = pto.bytewidth(idx_dtype)

    init_val = pto.PadValue.MIN.eval(src_dtype)

    if pto.constexpr(elem_bytes == 4):
        store_dist = pto.VStoreDist.ONE_POINT_B32
    elif pto.constexpr(elem_bytes == 2):
        store_dist = pto.VStoreDist.ONE_POINT_B16
    else:
        store_dist = pto.VStoreDist.ONE_POINT_B8

    for row in range(0, valid_rows, 1):
        remained = valid_cols

        v_val_acc = pto.vbr(init_val)
        init_zero_idx = idx_dtype(0)
        v_idx_acc = pto.vbr(init_zero_idx)

        mask_1, _ = pto.make_mask(src_dtype, 1)
        mask_1_idx, _ = pto.make_mask(idx_dtype, 1)

        for col in range(0, valid_cols, lanes):
            mask, remained = pto.make_mask(src_dtype, remained)
            v_src = pto.vlds(src[row, col:])
            v_reduced = pto.vcmax(v_src, mask)

            v_val, v_idx = pto.vdintlv(v_reduced, pto.vbr(src_dtype(0)))
            v_idx = pto.vbitcast(v_idx, idx_dtype)

            col_offset = idx_dtype(col)
            v_idx = pto.vadds(v_idx, col_offset, mask_1_idx)

            cmp_mask = pto.vcmp(v_val_acc, v_val, mask_1, "lt")

            v_val_acc = pto.vsel(v_val, v_val_acc, cmp_mask)
            cmp_mask_b32 = pto.pbitcast(cmp_mask, pto.mask_b32)
            v_idx_acc = pto.vsel(v_idx, v_idx_acc, cmp_mask_b32)

        pto.vsts(v_idx_acc, dst[row, 0:], mask_1_idx, dist=store_dist)
    return
\end{lstlisting}
\Needspace{0.34\textheight}
